\title{Torque-Referenced Consistency Regularization of dq Flux Maps}
\author{}
\date{7 October 2026}
\newcommand{\PaperAbstract}{Can an independently measured stationary torque expose and reduce spatially implausible structure in a reconstructed dq flux map? We model measured-minus-predicted torque as an unknown group offset plus spatial structure and noise; agreement therefore means a nearly flat residual within comparable speed, temperature and configuration, not pointwise equality. This yields a global minimum-change and smoothness-regularized flux correction. A scalar co-energy correction is used as the physically motivated numerical variant because it enforces reciprocity structurally, while remaining a prior rather than proof of magnetic truth. In a known-truth experiment it reduces flux RMSE and centered torque-residual spread; an injected radial potential remains invisible, as predicted by the retained rank-one observability theory. A portable two-system PSM dataset demonstrates offset removal and multi-speed transfer checks with $k_T=3p$. The CAN channel, shaft losses, terminal-current reconstruction and speed-dependent nuisance mechanisms remain experimentally unqualified, so no true-flux claim is made.}
\newif\ifshowglossary\showglossarytrue
\newif\ifshowbibliography\showbibliographytrue
